% Part: computability
% Chapter: recursive-functions
% Section: composition

\documentclass[../../../include/open-logic-section]{subfiles}

\begin{document}

\olfileid{cmp}{rec}{com}
\olsection{Komposisi}

Yen $f$ lan $g$ iku rong fungsi siji panggonan ing wilangan asli,
kita bisa nyusun komposisine: $h(x) = f(g(x))$. Fungsi anyar~$h(x)$
banjur ditegesi kanthi \emph{komposisi} saka fungsi $f$ lan~$g$.
Kita arep nggeneralisasi iki menyang fungsi kang nduweni argumen
luwih saka siji.

Mangkene salah siji carane: upamana $f$ iku fungsi $k$ panggonan,
lan $g_0$, \dots, $g_{k-1}$ iku $k$ fungsi kang kabeh nduweni
$n$ panggonan. Banjur kita bisa netepake fungsi anyar
kanthi $n$ panggonan, yaiku $h$, kaya mangkene:
\[
h(x_0, \dots, x_{n-1}) =
f(g_0(x_0, \dots, x_{n-1}), \dots, g_{k-1}(x_0, \dots, x_{n-1}))
\]
Yen $f$ lan kabeh $g_i$ komputabel, $h$ uga komputabel:
Kanggo ngitung $h(x_0,
\dots, x_{n-1})$, dhisik itungen biji $y_i = g_i(x_0, \dots,
x_{n-1})$ kanggo saben $i = 0$, \dots,~$k-1$. Banjur
lebokna biji kasebut menyang $f$ kanggo ngitung
% OLPL-122: Sumber nganggo indeks pungkasan k-1 kanggo input h;
% saka definisine, h duwe n input lan indeks pungkasan n-1.
$h(x_0, \dots, x_{n-1}) = f(y_0, \dots, y_{k-1})$. Cathetan koreksi sumber OLPL-122: input h nduweni n panggonan, dene input f nduweni k panggonan; sumber salah nggunakake k minus siji kanggo indeks pungkasan input h.

Iki bisa katon kaya gambaran kang banget mbatesi babagan
carane kita ngitung fungsi anyar nganggo fungsi kang wis ana.
Salah sijine, kadhangkala kita ora nggunakake kabeh argumen
sawijining fungsi, kayata nalika kita netepake
$g(x, y, z) = \Succ(z)$ kanggo definisi rekursif primitif
saka~$\Add$. Upamana kita bisa nggunakake fungsi-fungsi iki:
\[
\Proj{n}{i}(x_0, \dots, x_{n-1}) = x_i
\]
Fungsi~$\Proj{k}{i}$ diarani fungsi \emph{proyeksi}:
$\Proj{n}{i}$ iku fungsi $n$ panggonan. Cathetan panyunting: ing kabeh fungsi proyeksi iki, n positif lan i saka nol nganti n minus siji. Banjur $g$
bisa ditegesi kanthi
\[
g(x, y, z) = \Succ(\Proj{3}{2}(x, y, z)).
\]
Ing kene peran $f$ dimainake dening fungsi $1$ panggonan
$\Succ$, mula $k=1$. Lan kita nduweni fungsi $3$ panggonan
$\Proj{3}{2}$ kang dadi $g_0$. Asile yaiku fungsi $3$
panggonan kang ngasilake penerus argumen katelu.

Fungsi proyeksi uga ngidini kita netepake fungsi anyar
kanthi ngowahi urutan utawa nyamakake argumen. Umpamane,
fungsi $h(x)
= \Add(x, x)$ bisa ditegesi kanthi
\[
h(x_0) = \Add(\Proj{1}{0}(x_0),\Proj{1}{0}(x_0)).
\]
Ing kene $k=2$, $n=1$, peran $f(y_0,y_1)$ dimainake
dening $\Add$, lan peran $g_0(x_0)$ lan $g_1(x_0)$
kalorone dimainake dening~$\Proj{1}{0}(x_0)$,
yaiku fungsi proyeksi siji panggonan (uga diarani
fungsi identitas).

Yen $f(y_0, y_1)$ iku fungsi kang wis kita duweni,
kita bisa netepake fungsi $h(x_0, x_1) = f(x_1, x_0)$
kanthi
\[
h(x_0, x_1) = f(\Proj{2}{1}(x_0, x_1),\Proj{2}{0}(x_0, x_1)).
\]
Ing kene $k=2$, $n = 2$, lan peran $g_0$ lan $g_1$
dimainake dening $\Proj{2}{1}$ lan~$\Proj{2}{0}$
kanthi urutan kuwi.

Kita uga bisa kuwatir yen $g_0$, \dots,~$g_{k-1}$
kabeh kudu nduweni aritas~$n$. (Elinga yen \emph{aritas}
fungsi iku cacah argumene; fungsi $n$ panggonan nduweni
aritas~$n$.) Nanging tambahan fungsi proyeksi menehi
keluwesan kang dibutuhake. Umpamane, upamana $f$ lan~$g$
iku fungsi $3$ panggonan lan $h$~iku fungsi $2$ panggonan
kang ditegesi kanthi
\[
h(x,y) = f(x,g(x,x,y),y).
\]
Definisi~$h$ bisa ditulis maneh nganggo fungsi proyeksi,
yaiku
\[
h(x,y) = f(\Proj{2}{0}(x,y), g(\Proj{2}{0}(x,y), \Proj{2}{0}(x,y),
\Proj{2}{1}(x,y)), \Proj{2}{1}(x,y)).
\]
Banjur $h$ iku komposisi $f$ karo $\Proj{2}{0}$, $l$,
lan $\Proj{2}{1}$, ing ngendi
\[
l(x,y) = g(\Proj{2}{0}(x,y),\Proj{2}{0}(x,y),\Proj{2}{1}(x,y)),
\]
yaiku $l$ iku komposisi $g$ karo $\Proj{2}{0}$,
$\Proj{2}{0}$, lan~$\Proj{2}{1}$.

\end{document}
